\documentclass[a4paper,8pt,landscape]{article}

\usepackage[ngerman]{babel}
\usepackage[T1]{fontenc}
\usepackage[utf8]{inputenc}
\usepackage[a4paper,margin=0.9cm]{geometry}
\usepackage{amsmath,amssymb}
\usepackage{multicol}
\usepackage{enumitem}
\usepackage[hidelinks]{hyperref}

\setlength{\columnsep}{0.45cm}
\setlength{\parindent}{0pt}
\setlength{\parskip}{0.15em}
\setlist[itemize]{leftmargin=1.2em,itemsep=0.05em,topsep=0.05em}
\setlist[enumerate]{leftmargin=1.4em,itemsep=0.05em,topsep=0.05em}

\newcommand{\dd}{\,\mathrm{d}}
\newcommand{\Var}{\mathrm{Var}}
\newcommand{\Cov}{\mathrm{Cov}}
\newcommand{\Corr}{\mathrm{Corr}}
\newcommand{\Bin}{\mathrm{Bin}}
\newcommand{\Geom}{\mathrm{Geom}}
\newcommand{\Hyp}{\mathrm{Hyp}}
\newcommand{\Exp}{\mathrm{Exp}}
\newcommand{\Unif}{\mathrm{Unif}}
\newcommand{\Normal}{\mathcal{N}}
\newcommand{\Ber}{\mathrm{Ber}}
\newcommand{\EW}{\mathrm{E}}
\newcommand{\Pp}{\mathrm{P}}

\begin{document}
\scriptsize

\begin{center}
{\Large Stochastik Cheatsheet -- Kompaktfassung}\\
Formeln, Notation, Einsatzregeln
\end{center}
\vspace{-0.4em}

\begin{multicols*}{3}

\section*{1. Grundnotation}
\[
\Omega:\text{ Ergebnismenge},\quad \omega\in\Omega,\quad A,B\subseteq\Omega
\]
\[
\overline A=\Omega\setminus A,\quad A\cup B,\quad A\cap B,\quad A\setminus B
\]
\[
P(\emptyset)=0,\quad P(\Omega)=1,\quad P(\overline A)=1-P(A)
\]
\[
P(A\cup B)=P(A)+P(B)-P(A\cap B)
\]
\[
A\cap B=\emptyset \Rightarrow P(A\cup B)=P(A)+P(B)
\]
\[
P(A\setminus B)=P(A)-P(A\cap B)
\]
\[
P(B\setminus A)=P(B)-P(A\cap B)
\]
\[
P(\overline A\cap B)=P(B)-P(A\cap B)
\]
\[
P(\overline A\cup \overline B)=1-P(A\cap B)
\]
de Morgan:
\[
\overline{A\cup B}=\overline A\cap \overline B
\]
\[
\overline{A\cap B}=\overline A\cup \overline B
\]
Laplace:
\[
P(A)=\frac{|A|}{|\Omega|}
\]
wenn alle Elementarergebnisse gleich wahrscheinlich sind.

\textbf{Notation in Loesungen}
\begin{itemize}
  \item ``Sei $\Omega=\{\dots\}$ die Ergebnismenge.''
  \item ``Sei $A=\{\dots\}$ das Ereignis \dots''
  \item ``Da Laplace-Experiment / da $A\cap B=\emptyset$ / mit Bayes folgt \dots''
\end{itemize}

\textbf{Merke}
\begin{itemize}
  \item erst $\Omega$, dann Ereignisse, dann Formel
  \item ``oder'' $=$ Vereinigung, ``und'' $=$ Schnitt
  \item Gegenereignis oft kuerzer
\end{itemize}

\section*{2. Wahrscheinlichkeitsfunktion und ZV}
Diskret:
\[
p(\omega)=P(\{\omega\}),\qquad p(\omega)\ge 0,\qquad \sum_{\omega\in\Omega}p(\omega)=1
\]
\[
P(A)=\sum_{\omega\in A}p(\omega)
\]
Wenn $p(i)=c f(i)$:
\[
\sum_i c f(i)=1 \Rightarrow c=\frac{1}{\sum_i f(i)}
\]

\textbf{Template fuer Wahrscheinlichkeitsfunktionen}
\begin{itemize}
  \item 1. Ergebnismenge festlegen: $\Omega=\{\omega_1,\dots,\omega_n\}$
  \item 2. Jedem Ergebnis eine Einzelwahrscheinlichkeit geben: $p(\omega_i)$
  \item 3. Pruefen: $p(\omega_i)\ge 0$ und $\sum_{\omega\in\Omega}p(\omega)=1$
  \item 4. Ereigniswahrscheinlichkeit ueber Summe: $P(A)=\sum_{\omega\in A}p(\omega)$
\end{itemize}

\textbf{Wuerfelbeispiel}
\[
\Omega=\{1,2,3,4,5,6\},\qquad p(i)=\frac16 \text{ fuer } i=1,\dots,6
\]
\[
\sum_{i=1}^6 p(i)=6\cdot \frac16=1
\]
Fuer $A=\{2,4,6\}$:
\[
P(A)=p(2)+p(4)+p(6)=\frac16+\frac16+\frac16=\frac12
\]

\textbf{Wenn Konstante gesucht ist}
\[
p(i)=c\cdot f(i)\ \Rightarrow\ \sum_i p(i)=1\ \Rightarrow\ c=\frac{1}{\sum_i f(i)}
\]

\textbf{Sauberer Klausursatz}
\begin{itemize}
  \item ``Definiere $p:\Omega\to[0,1]$ durch $p(i)=\dots$.''
  \item ``Da $p(i)\ge 0$ und $\sum_{\omega\in\Omega}p(\omega)=1$, ist $p$ eine Wahrscheinlichkeitsfunktion.''
\end{itemize}

Zufallsvariable:
\[
X:\Omega\to\mathbb{R},\qquad P_X(B)=P(X\in B)=P(X^{-1}(B))
\]
Diskret:
\[
P(X=x)=\sum_{\omega:X(\omega)=x}P(\{\omega\})
\]

\textbf{Typische Fallen}
\begin{itemize}
  \item Verteilung von $X$ ist oft \emph{nicht} gleichverteilt, auch wenn $\Omega$ gleichverteilt ist.
  \item Nicht unterscheidbare Objekte nie blind wie verschiedene zaehlen.
\end{itemize}

\section*{3. Kombinatorik}
\textbf{Schublade}
\begin{center}
\begin{tabular}{ll}
Reihenfolge wichtig? & mit/ohne Wiederholung?
\end{tabular}
\end{center}

Permutation ohne Wdh.:
\[
n!
\]
Permutation mit Wdh. bei Multimenge $(n_1,\dots,n_r)$:
\[
\frac{n!}{n_1!\cdots n_r!}
\]
Variation ohne Wdh.:
\[
V(n,k)=\frac{n!}{(n-k)!}
\]
Variation mit Wdh.:
\[
n^k
\]
Kombination ohne Wdh.:
\[
\binom{n}{k}=\frac{n!}{k!(n-k)!}
\]
Kombination mit Wdh.:
\[
\binom{n+k-1}{k}
\]

\textbf{Standardfaelle}
\begin{itemize}
  \item Lotto, Teams, Auswahlmenge $\Rightarrow \binom{n}{k}$
  \item PIN, Code, geordnete Ziehung mit Zuruecklegen $\Rightarrow n^k$
  \item Sitzordnung / Reihenfolge $\Rightarrow n!$ oder $\frac{n!}{(n-k)!}$
\end{itemize}

\textbf{Eselsbruecke}
\begin{itemize}
  \item Permutation: alle $n$ Objekte werden angeordnet
  \item Variation: nur $k$ aus $n$, Reihenfolge wichtig
  \item Kombination: nur $k$ aus $n$, Reihenfolge egal
\end{itemize}

Ein-/Ausschluss:
\[
|A\cup B|=|A|+|B|-|A\cap B|
\]
\[
P(A\cup B)=P(A)+P(B)-P(A\cap B)
\]

Geburtstagsparadoxon:
\[
P(\text{mind. 2 gleich})=1-P(\text{alle verschieden})
\]
\[
P(\text{alle verschieden})=\frac{365\cdot364\cdots(365-n+1)}{365^n}
\]

\section*{4. Bedingte Wkt., Bayes, Unabhaengigkeit}
Definition:
\[
P(A\mid B)=\frac{P(A\cap B)}{P(B)} \quad (P(B)>0)
\]
Umformung:
\[
P(A\cap B)=P(A\mid B)P(B)=P(B\mid A)P(A)
\]
Totale Wahrscheinlichkeit bei Zerlegung $B_1,\dots,B_n$:
\[
P(A)=\sum_{i=1}^n P(A\mid B_i)P(B_i)
\]
Bayes:
\[
P(B_j\mid A)=\frac{P(A\mid B_j)P(B_j)}{\sum_{i=1}^n P(A\mid B_i)P(B_i)}
\]
Unabhaengigkeit:
\[
A\perp B \iff P(A\cap B)=P(A)P(B)
\]
Dann auch:
\[
P(A\mid B)=P(A),\qquad P(B\mid A)=P(B)
\]

\textbf{Wann?}
\begin{itemize}
  \item vorwaerts durch Baum $\Rightarrow$ Produkt + totale Wkt.
  \item rueckwaerts nach Beobachtung $\Rightarrow$ Bayes
  \item ``unabhaengig?'' $\Rightarrow$ Produktkriterium pruefen
\end{itemize}

\textbf{Falle}
\begin{itemize}
  \item disjunkt $\neq$ unabhaengig
  \item $P(A\mid B)\neq P(B\mid A)$
\end{itemize}

\section*{5. Diskrete Standardverteilungen}
\textbf{Merkschema}
\begin{itemize}
  \item Einsatz: Wann passt die Verteilung?
  \item Wertebereich $W(X)$: Welche Werte kann $X$ annehmen?
  \item Formel: $P(X=k)$
  \item Lage/Streuung: $E(X)$ und $\Var(X)$
\end{itemize}

Bernoulli:
\[
X\sim\Ber(p),\quad P(X=1)=p,\quad P(X=0)=1-p
\]
\[
E(X)=p,\quad \Var(X)=p(1-p)
\]

Binomial:
\[
X\sim\Bin(n,p),\quad P(X=k)=\binom{n}{k}p^k(1-p)^{n-k}
\]
\[
E(X)=np,\quad \Var(X)=np(1-p)
\]
Nutzen: $n$ feste unabhaengige Versuche, Trefferzahl.

Geometrisch:
\[
X\sim\Geom(p),\quad P(X=k)=(1-p)^{k-1}p,\quad k\ge 1
\]
\[
P(X>m)=(1-p)^m
\]
\[
E(X)=\frac{1}{p},\quad \Var(X)=\frac{1-p}{p^2}
\]
Nutzen: Warten bis zum ersten Treffer.

Hypergeometrisch:
\[
X\sim\Hyp(N,M,n),\quad P(X=k)=\frac{\binom{M}{k}\binom{N-M}{n-k}}{\binom{N}{n}}
\]
\[
E(X)=n\frac{M}{N},\quad \Var(X)=n\frac{M}{N}\left(1-\frac{M}{N}\right)\frac{N-n}{N-1}
\]
Nutzen: Ziehen ohne Zuruecklegen.

Poisson:
\[
X\sim\mathrm{Poi}(\lambda),\quad W(X)=\{0,1,2,\dots\}
\]
\[
P(X=k)=\frac{\lambda^k e^{-\lambda}}{k!}
\]
\[
E(X)=\lambda,\quad \Var(X)=\lambda
\]
Nutzen: Anzahl seltener Ereignisse pro Zeit-/Raumintervall.

\textbf{Entscheidung}
\begin{itemize}
  \item mit Zuruecklegen / unabhaengig $\Rightarrow$ binomial
  \item bis erster Treffer $\Rightarrow$ geometrisch
  \item ohne Zuruecklegen $\Rightarrow$ hypergeometrisch
  \item seltene Ereignisse pro Intervall $\Rightarrow$ Poisson
\end{itemize}

\section*{6. Erwartungswert und Varianz}
Diskret:
\[
E(X)=\sum_x x\,P(X=x)
\]
oder direkt ueber $\Omega$:
\[
E(X)=\sum_{\omega\in\Omega}X(\omega)P(\{\omega\})
\]
Allgemein:
\[
E(g(X))=\sum_x g(x)P(X=x)
\]

Linearitaet:
\[
E(aX+bY+c)=aE(X)+bE(Y)+c
\]
gilt auch bei \textbf{abhaengigen} Variablen.
Bei Unabhaengigkeit:
\[
E(XY)=E(X)E(Y)
\]

Varianz:
\[
\Var(X)=E\big((X-E(X))^2\big)=E(X^2)-E(X)^2
\]
\[
\sigma_X=\sqrt{\Var(X)}
\]
Rechenregeln:
\[
\Var(aX+b)=a^2\Var(X)
\]
\[
\Var(X+Y)=\Var(X)+\Var(Y)+2\Cov(X,Y)
\]
bei Unabhaengigkeit:
\[
\Var(X+Y)=\Var(X)+\Var(Y)
\]

\textbf{Wichtig}
\begin{itemize}
  \item Fuer $E(X+Y)$ brauchst du \emph{keine} Unabhaengigkeit:
  \[
  E(X+Y)=E(X)+E(Y)
  \]
  \item Abhaengigkeit ist erst relevant bei $E(XY)$, Varianz von Summen und Kovarianz.
\end{itemize}

Tschebyscheff:
\[
P(|X-E(X)|\ge \varepsilon)\le \frac{\Var(X)}{\varepsilon^2}
\]
\[
P(|X-E(X)|<\varepsilon)\ge 1-\frac{\Var(X)}{\varepsilon^2}
\]

\textbf{Taktik}
\begin{itemize}
  \item Gewinnfunktion erst als $g(X)$ formulieren
  \item dann $E(g(X))$ statt Fall-fuer-Fall-Geschwafel
\end{itemize}

\section*{7. Kovarianz, Korrelation, LLN}
\[
\Cov(X,Y)=E(XY)-E(X)E(Y)
\]
\[
\Corr(X,Y)=\frac{\Cov(X,Y)}{\sigma_X\sigma_Y}
\]
\[
-1\le \Corr(X,Y)\le 1
\]
Unabhaengigkeit $\Rightarrow \Cov(X,Y)=0$, aber nicht umgekehrt.

Fuer Summe:
\[
\Var\Big(\sum_{i=1}^n X_i\Big)=\sum_{i=1}^n \Var(X_i)+2\sum_{i<j}\Cov(X_i,X_j)
\]
bei paarweiser Unabhaengigkeit:
\[
\Var\Big(\sum_{i=1}^n X_i\Big)=\sum_{i=1}^n \Var(X_i)
\]

Mittelwert:
\[
\overline X_n=\frac{1}{n}\sum_{i=1}^n X_i
\]
iid mit $E(X_i)=\mu$, $\Var(X_i)=\sigma^2$:
\[
E(\overline X_n)=\mu,\qquad \Var(\overline X_n)=\frac{\sigma^2}{n}
\]
Schwaches Gesetz:
\[
P(|\overline X_n-\mu|\ge \varepsilon)\le \frac{\sigma^2}{n\varepsilon^2}\to 0
\]
Relative Haeufigkeit als Spezialfall:
\[
\hat p_n=\frac{X_1+\cdots+X_n}{n},\quad X_i\sim\Ber(p)
\]
\[
E(\hat p_n)=p,\qquad \Var(\hat p_n)=\frac{p(1-p)}{n}
\]

\section*{8. Allgemein und stetig}
Wahrscheinlichkeitsraum allgemein:
\[
(\Omega,\mathcal{A},P)
\]
$\mathcal{A}$: Sigma-Algebra.

Verteilungsfunktion:
\[
F_X(x)=P(X\le x)
\]
\[
P(a<X\le b)=F_X(b)-F_X(a)
\]
Eigenschaften: monoton, rechtsstetig, Grenzwerte $0/1$.

Stetige ZV mit Dichte $f$:
\[
f(x)\ge 0,\qquad \int_{-\infty}^{\infty}f(x)\dd x=1
\]
\[
P(a\le X\le b)=\int_a^b f(x)\dd x
\]
\[
F(x)=\int_{-\infty}^x f(t)\dd t,\qquad f=F'
\]
\[
P(X=x)=0
\]

Erwartungswert / Varianz stetig:
\[
E(X)=\int_{-\infty}^{\infty}x f(x)\dd x
\]
\[
E(g(X))=\int_{-\infty}^{\infty}g(x)f(x)\dd x
\]
\[
\Var(X)=E(X^2)-E(X)^2
\]

\section*{9. Stetige Standardverteilungen}
\textbf{Merkschema}
\begin{itemize}
  \item Einsatz: Welche Situation beschreibt $X$?
  \item Dichte $f(x)$ und oft auch Verteilungsfunktion $F(x)$
  \item Lage/Streuung: $E(X)$ und $\Var(X)$
\end{itemize}

Gleichverteilung auf $[a,b]$:
\[
X\sim \Unif[a,b],\quad f(x)=\frac{1}{b-a}
\]
\[
E(X)=\frac{a+b}{2},\qquad \Var(X)=\frac{(b-a)^2}{12}
\]

Exponentialverteilung:
\[
X\sim \Exp(\lambda),\quad f(x)=\lambda e^{-\lambda x}\ \ (x\ge 0)
\]
\[
F(x)=1-e^{-\lambda x},\quad E(X)=\frac{1}{\lambda},\quad \Var(X)=\frac{1}{\lambda^2}
\]
Gedaechtnislos:
\[
P(X>s+t\mid X>s)=P(X>t)
\]

Normalverteilung:
\[
X\sim \Normal(\mu,\sigma^2),\quad f(x)=\frac{1}{\sqrt{2\pi}\sigma}e^{-(x-\mu)^2/(2\sigma^2)}
\]
Standardisierung:
\[
Z=\frac{X-\mu}{\sigma}\sim \Normal(0,1)
\]
\[
P(X\le x)=\Phi\left(\frac{x-\mu}{\sigma}\right)
\]
Symmetrie:
\[
\Phi(-z)=1-\Phi(z)
\]
\[
P(|Z|\le z)=2\Phi(z)-1
\]
$3\sigma$-Regel:
\[
P(|X-\mu|\le \sigma)\approx 0.683,\ 
P(|X-\mu|\le 2\sigma)\approx 0.955,\ 
P(|X-\mu|\le 3\sigma)\approx 0.997
\]

Linearkombinationen:
\[
aX+b\sim \Normal(a\mu+b,a^2\sigma^2)
\]
Bei unabhaengigen Normalvariablen:
\[
X+Y\sim \Normal(\mu_X+\mu_Y,\sigma_X^2+\sigma_Y^2)
\]
\[
X-Y\sim \Normal(\mu_X-\mu_Y,\sigma_X^2+\sigma_Y^2)
\]

\section*{10. Binomial $\approx$ Normal}
Wenn $X\sim\Bin(n,p)$ und $np(1-p)$ gross genug:
\[
X\approx \Normal(np,np(1-p))
\]
Standardisiert:
\[
Z=\frac{X-np}{\sqrt{np(1-p)}}\approx \Normal(0,1)
\]
Stetigkeitskorrektur:
\[
P(X\le k)\approx P\left(Y\le k+\frac12\right)
\]
\[
P(X\ge k)\approx P\left(Y\ge k-\frac12\right)
\]
\[
P(a\le X\le b)\approx P\left(a-\frac12\le Y\le b+\frac12\right)
\]
mit $Y\sim\Normal(np,np(1-p))$.

\textbf{Merkkasten Stetigkeitskorrektur}
\begin{itemize}
  \item diskret $\to$ stetig: Grenze um $0.5$ verschieben
  \item $P(X\le k)\Rightarrow k+0.5$
  \item $P(X\ge k)\Rightarrow k-0.5$
  \item $P(X=k)\Rightarrow k-0.5 \le Y \le k+0.5$
  \item Intervall $[a,b]\Rightarrow [a-0.5,b+0.5]$
\end{itemize}

\section*{11. Beschreibende Statistik}
Stichprobe:
\[
x_1,\dots,x_n
\]
Relative Haeufigkeit:
\[
\hat f(a)=\frac{1}{n}\#\{i:x_i=a\}
\]
Empirische Verteilungsfunktion:
\[
\hat F_n(x)=\frac{1}{n}\#\{i:x_i\le x\}
\]
Arithmetisches Mittel:
\[
\overline x=\frac{1}{n}\sum_{i=1}^n x_i
\]
Median: mittlerer geordneter Wert.

Stichprobenvarianz:
\[
s_n^2=\frac{1}{n}\sum_{i=1}^n (x_i-\overline x)^2
\]
Unverzerrte Varianz:
\[
s^2=\frac{1}{n-1}\sum_{i=1}^n (x_i-\overline x)^2
\]
\[
s=\sqrt{s^2}
\]

\textbf{Merke}
\begin{itemize}
  \item $\hat f,\hat F_n$: Datenbeschreibung
  \item $\bar x$: Lage
  \item $s^2,s$: Streuung
\end{itemize}

\section*{12. Schaetzen}
Parameter $\theta$, Schaetzer $T=T(X_1,\dots,X_n)$.

Punktschaetzer:
\[
\hat p=\overline X\quad \text{bei } X_i\sim\Ber(p)
\]
\[
\hat \mu=\overline X
\]
\[
\hat \sigma^2=\frac{1}{n-1}\sum_{i=1}^n (X_i-\overline X)^2
\]

Wichtige Eigenschaften:
\begin{itemize}
  \item erwartungstreu: $E(T)=\theta$
  \item konsistent: $T_n\to \theta$
  \item effizient: kleinere Varianz als andere erwartungstreue Schaetzer
\end{itemize}

\textbf{Klassiker}
\begin{itemize}
  \item Anteil schaetzen $\Rightarrow \hat p=\overline X$
  \item Mittelwert schaetzen $\Rightarrow \bar X$
  \item Varianz aus Daten $\Rightarrow \frac{1}{n-1}\sum (X_i-\bar X)^2$
\end{itemize}

\section*{13. Hypothesentests}
Hypothesen:
\[
H_0:\theta=\theta_0 \quad\text{gegen}\quad H_1:\theta\neq \theta_0,\ >\theta_0,\ <\theta_0
\]
Fehler 1. Art:
\[
\alpha=P(\text{$H_0$ verwerfen}\mid H_0\text{ wahr})
\]
Fehler 2. Art:
\[
\beta(\theta)=P(\text{$H_0$ beibehalten}\mid \theta\in H_1)
\]
Guetefunktion:
\[
g(\theta)=P_\theta(\text{$H_0$ verwerfen})=1-\beta(\theta)
\]

Zweiseitiger Test fuer Mittelwert bei bekannter $\sigma$:
\[
Z=\frac{\overline X-\mu_0}{\sigma/\sqrt n}
\]
Unter $H_0$:
\[
Z\sim \Normal(0,1)
\]
Ablehnen von $H_0$ zum Niveau $\alpha$, wenn:
\[
|Z|>z_{1-\alpha/2}
\]
Typische Quantile:
\[
z_{0.95}\approx 1.645,\quad z_{0.975}\approx 1.96,\quad z_{0.995}\approx 2.576
\]

\textbf{Ablauf}
\begin{enumerate}
  \item $H_0,H_1$ sauber aufschreiben
  \item Teststatistik angeben
  \item Verteilung unter $H_0$
  \item kritischen Bereich / p-Wert
  \item Entscheidung im Kontext
\end{enumerate}

\section*{14. Mini-Entscheidungsbaum}
\begin{itemize}
  \item Anzahl guenstiger Faelle? $\Rightarrow$ Kombinatorik + Laplace
  \item Wahrscheinlichkeiten nach Beobachtung? $\Rightarrow$ Bayes
  \item Trefferzahl in $n$ Versuchen? $\Rightarrow$ Binomial
  \item Wartezeit bis erster Treffer? $\Rightarrow$ Geometrisch
  \item Ziehen ohne Zuruecklegen? $\Rightarrow$ Hypergeometrisch
  \item Mittelwert / Fairness / Gewinn? $\Rightarrow$ Erwartungswert
  \item Streuung / Fehlerband? $\Rightarrow$ Varianz, $\sigma$, Tschebyscheff
  \item Stetige Dichte gegeben? $\Rightarrow$ integrieren
  \item Normal verteilt / Approximation? $\Rightarrow$ standardisieren
  \item Daten gegeben? $\Rightarrow$ $\bar x$, Median, $s^2$, $\hat F_n$
  \item Parameter aus Stichprobe? $\Rightarrow$ Schaetzer
  \item Behauptung pruefen? $\Rightarrow$ Test
\end{itemize}

\section*{15. Typische Klausurmuster}
\textbf{ZV als Abbildung definieren}
\begin{itemize}
  \item ``$X:\Omega\to\mathbb{R},\ X(\omega)=$ Anzahl / Laenge / Trefferzahl \dots''
  \item Bildmenge explizit angeben, z.B. $W(X)=\{0,\dots,n\}$
\end{itemize}

\textbf{Gemeinsame Verteilung / Unabhaengigkeit}
\begin{itemize}
  \item Tabelle fuellen: alle moeglichen Paare $(x,y)$ systematisch zaehlen
  \item unabhaengig? $\Rightarrow P(X=x,Y=y)\stackrel{?}=P(X=x)P(Y=y)$
  \item oder bei Ereignissen: $P(A\cap B)\stackrel{?}=P(A)P(B)$
\end{itemize}

\textbf{Kovarianz / Korrelation}
\[
\Cov(X,Y)=\rho(X,Y)\sigma_X\sigma_Y
\]
\[
\rho(X,Y)=\frac{\Cov(X,Y)}{\sigma_X\sigma_Y}
\]
\[
\Var(X+Y)=\Var(X)+\Var(Y)+2\Cov(X,Y)
\]

\textbf{Anteil schaetzen / Konfidenzintervall}
\[
\hat p=r_n=\frac{1}{n}\sum_{i=1}^n X_i
\]
fuer grosses $n$:
\[
\hat p \pm z_{(1+\gamma)/2}\sqrt{\frac{\hat p(1-\hat p)}{n}}
\]

\textbf{Konfidenzintervall fuer $\mu$ bei bekannter $\sigma$}
\[
\overline X \pm z_{(1+\gamma)/2}\frac{\sigma}{\sqrt n}
\]

\textbf{Tschebyscheff fuer relative Haeufigkeit}
\[
P(|r_n-p|\ge \varepsilon)\le \frac{p(1-p)}{n\varepsilon^2}\le \frac{1}{4n\varepsilon^2}
\]
also fuer Schranke $\le \delta$:
\[
n\ge \frac{1}{4\delta\varepsilon^2}
\]

\textbf{Ungleichungen umformen}
\begin{itemize}
  \item aus $a/b \le c$ mit $b>0$ folgt $a \le bc$
  \item aus $a/b \ge c$ mit $b>0$ folgt $a \ge bc$
  \item mit negativer Zahl multiplizieren/dividieren $\Rightarrow$ Ungleichheitszeichen dreht sich um
  \item bei $P(|X-\mu|<\varepsilon)$ oft erst Mittelpunkt erkennen
  \item Intervallform: $a<X<b \iff |X-\frac{a+b}{2}|<\frac{b-a}{2}$
  \item analog: $a\le X\le b \iff |X-\frac{a+b}{2}|\le \frac{b-a}{2}$
\end{itemize}

\textbf{Typische Klausur-Umformungen}
\begin{itemize}
  \item $P(220<Z<260)=P(|Z-240|<20)$
  \item $P(|r_n-p|>0.1)\le 0.05 \Rightarrow \frac{1}{4n\cdot 0.1^2}\le 0.05 \Rightarrow n\ge 500$
  \item $P(Z>m)\le 0.05 \iff P(Z\le m)\ge 0.95$
  \item ``mindestens 90\%'' $\Rightarrow P(\dots)\ge 0.9$
  \item ``hoechstens 5\%'' $\Rightarrow P(\dots)\le 0.05$
\end{itemize}

\textbf{Normalapproximation der Binomialverteilung}
\begin{itemize}
  \item erst $X\sim\Bin(n,p)$ erkennen
  \item dann $X\approx \Normal(np,np(1-p))$
  \item Stetigkeitskorrektur nicht vergessen
\end{itemize}

\textbf{Einseitiger Test fuer Anteil oder Mittelwert}
\begin{itemize}
  \item ``haeufiger / groesser'' $\Rightarrow H_1:\theta>\theta_0$
  \item ``seltener / kleiner'' $\Rightarrow H_1:\theta<\theta_0$
  \item ``anders'' $\Rightarrow H_1:\theta\neq\theta_0$
\end{itemize}

\textbf{Stetige Dichteaufgaben}
\begin{itemize}
  \item Konstante $c$ suchen: $\int f(x)\dd x=1$
  \item dann $E(X)=\int x f(x)\dd x$, $\Var(X)=E(X^2)-E(X)^2$
\end{itemize}

\textbf{Keywords $\Rightarrow$ Methode}
\begin{itemize}
  \item ``genau $k$ Treffer'' $\Rightarrow$ meist Binomial: $\binom{n}{k}p^k(1-p)^{n-k}$
  \item ``mindestens / hoechstens'' $\Rightarrow$ Binomialsumme oder Gegenereignis
  \item ``bis zum ersten Treffer'' / ``erst beim $k$-ten Mal'' $\Rightarrow$ geometrisch
  \item ``ohne Zuruecklegen'' $\Rightarrow$ hypergeometrisch
  \item ``mit Zuruecklegen'' / ``unabhaengige Wiederholungen'' $\Rightarrow$ binomial
  \item ``in wie vielen Anordnungen'' $\Rightarrow$ Kombinatorik
  \item ``gleiche Objekte nicht unterscheiden'' $\Rightarrow$ Permutation mit Wiederholung
  \item ``gemeinsame Verteilung'' / Tabelle $\Rightarrow$ alle Paare $(x,y)$ systematisch bestimmen
  \item ``stochastisch unabhaengig?'' $\Rightarrow$ Produktkriterium
  \item ``im Mittel'' $\Rightarrow$ Erwartungswert
  \item ``Streuung'' / ``Standardabweichung'' $\Rightarrow$ Varianz, $\sigma$
  \item ``untere Schranke fuer Wahrscheinlichkeit'' $\Rightarrow$ oft Tschebyscheff
  \item ``mit hoechstens 5\% Wahrscheinlichkeit ausverkauft / nicht ausreichen'' $\Rightarrow$ Quantil, oft Normalapproximation
  \item ``approximiere durch Normalverteilung'' $\Rightarrow$ de Moivre-Laplace / CLT + Stetigkeitskorrektur
  \item ``schaetzen'' / ``Anteil in Grundgesamtheit'' $\Rightarrow$ $\hat p=r_n$
  \item ``Konfidenzintervall'' $\Rightarrow$ Schaetzer $\pm$ Quantil $\cdot$ Standardfehler
  \item ``Hypothesentest'' / ``Verdacht'' / ``behauptet'' $\Rightarrow$ $H_0,H_1$, Teststatistik, kritischer Bereich
  \item ``gegeben, dass'' $\Rightarrow$ bedingte Wahrscheinlichkeit
  \item ``Ursache nach Beobachtung'' $\Rightarrow$ Bayes
  \item ``Dichtefunktion'' $\Rightarrow$ integrieren, nie summieren
  \item ``Parameter $c$ bestimmen'' bei Dichte/Wkt.-Funktion $\Rightarrow$ Normierung auf $1$
  \item ``normalverteilt'' $\Rightarrow$ standardisieren: $Z=\frac{X-\mu}{\sigma}$
\end{itemize}

\textbf{Typische Formulierungen sauber uebersetzen}
\begin{itemize}
  \item ``mindestens $k$'' $\Rightarrow P(X\ge k)$
  \item ``mehr als $k$'' $\Rightarrow P(X>k)$
  \item ``hoechstens $k$'' $\Rightarrow P(X\le k)$
  \item ``weniger als $k$'' $\Rightarrow P(X<k)$
  \item ``genau $k$'' $\Rightarrow P(X=k)$
  \item ``zwischen $a$ und $b$'' $\Rightarrow P(a\le X\le b)$ oder Aufgabe genau lesen
\end{itemize}

\section*{16. Reicht das?}
\begin{itemize}
  \item Fuer die bisher gesehenen Klausurtypen: ziemlich weit ja.
  \item Nicht garantiert: jede Spezialfrage oder jede Theorie-/Begruendungsfrage.
  \item Wenn du haengst: erst Keyword suchen, dann passende Schublade, dann Formel.
\end{itemize}

\section*{17. Notation, die Punkte bringt}
\begin{itemize}
  \item Ereignisse als Mengen notieren, nicht nur in Worten
  \item Zufallsvariable benennen: ``Sei $X$ = Anzahl der Treffer''
  \item Verteilung sauber: ``$X\sim\Bin(n,p)$''
  \item Formeln mit Begruendung: ``da unabhaengig'', ``da ohne Zuruecklegen'', ``mit Bayes''
  \item am Ende Antwortsatz mit Kontext
\end{itemize}

\section*{18. Symbollegende}
\begin{itemize}
  \item $\omega$: einzelnes Ergebnis eines Zufallsexperiments
  \item $\Omega$: Ergebnismenge, also alle moeglichen Ergebnisse
  \item $A,B$: Ereignisse, also Teilmengen von $\Omega$
  \item $\emptyset$: unmoegliches Ereignis, leere Menge
  \item $\overline A$: Gegenereignis zu $A$
  \item $P(A)$: Wahrscheinlichkeit des Ereignisses $A$
  \item $p$: Trefferwahrscheinlichkeit bei Bernoulli/Binomial/Geometrisch
  \item $q=1-p$: Gegenwahrscheinlichkeit
  \item $X,Y$: Zufallsvariablen
  \item $x_i$: einzelne Beobachtung / einzelner Datenwert
  \item $n$: Stichprobenumfang oder Anzahl der Versuche
  \item $k$: konkrete Trefferzahl / Auspraegung
  \item $\mu$: Erwartungswert / Mittelwert in der Grundgesamtheit
  \item $\mu_0$: Wert aus der Nullhypothese
  \item $\sigma$: Standardabweichung in der Grundgesamtheit
  \item $\sigma^2$: Varianz in der Grundgesamtheit
  \item $\overline X,\overline x$: Stichprobenmittel
  \item $s^2,s$: Stichprobenvarianz und Stichproben-Standardabweichung
  \item $\theta$: allgemeiner unbekannter Parameter
  \item $\hat \theta$: Schaetzer fuer $\theta$
  \item $\hat p$: geschaetzter Anteil, meist $\overline X$
  \item $\lambda$: Rate der Exponentialverteilung
  \item $\rho$: Korrelationskoeffizient, also standardisierte Abhaengigkeit
  \item $\Cov(X,Y)$: Kovarianz von $X$ und $Y$
  \item $\Var(X)$: Varianz von $X$
  \item $\Phi(z)$: Verteilungsfunktion der Standardnormalverteilung
  \item $z_{1-\alpha/2}$: Quantil der Standardnormalverteilung
  \item $\alpha$: Fehler 1. Art / Signifikanzniveau
  \item $\beta$: Fehler 2. Art
  \item $H_0,H_1$: Nullhypothese und Alternative
\end{itemize}

\end{multicols*}
\end{document}
